\section{Conclusion}\label{sec:conclusion}

The shortfall of the Federal Reserve's mortgage runoff beneath its Quantitative Tightening caps was not solely a product of high interest rates. It followed from the structural design of the U.S. mortgage contract interacting with a prepayment environment mostly rate-inelastic relative to the caps. Under the production floor form the $\beta_1 = 0$ null recovers 85.7\% of the benchmark at the headline calibration (35.6\% under the additive form) and 88.7\% at the in-sample calibration point. So the caps could not have bound under any rate response above the baseline turnover the null embeds. What delivers the level is scheduled amortization, the seasoning baseline, and the baseline turnover floor. The paper's own ablations bound the named loan-level mechanisms against that base: burnout $0.45$ points at the headline floor and $0.86$ in-sample, FICO/LTV covariates $1.3$ points, the delinquency pipeline \$0.39 billion. No servicer channel beyond that pipeline is modeled at all.

That decomposition is established in Section~\ref{sec:hazard} on real Freddie Mac loan-level data, in a framework built entirely from loan-level survival dynamics with no household utility function at all. In its microsimulation form, on the benchmark-consistent accounting basis, it recovers 91.3\% of the full empirical benchmark at the off-window floor this paper headlines, 88.2--93.7\% across that floor's defensible range, and 97.9\% at the in-sample calibration point. So the miss is 8.7 points on the headline calibration and 2.1 points on the in-sample one. That is a level-accuracy distinction, not a timing one: no estimator's monthly co-movement survives detrending (Section~\ref{sec:pathb}). The agent-based model of Section~\ref{sec:abm}, which runs on the synthetic household population specified in Section~\ref{sec:method}, stands beside that finding as a stress-test companion---rational, behaviorally extended, or subjected to a direct burnout test, it cannot explain the majority of the Federal Reserve's cap shortfall. Neither is derived from the other, but they are not independent in calibration: both are anchored to a baseline turnover floor read off discount-cohort turnover, and at the in-sample calibration point to the same 2023--24 read (Section~\ref{sec:pathb}).

The lock-in marginal is the secondary result, and it is reported as a bound. Against the Federal Reserve's own ex-ante projection rather than the non-binding cap, the lock-in channel accounts for roughly half the genuine surprise under the central allocation; that share runs 23--49\% at the off-window floor and 38--80\% at the in-sample calibration point across the two disclosed intra-2022 allocations (Section~\ref{sec:method-benchmark}). The elasticity's contribution is a bounded addition over that rate-inelastic baseline: $+2.3$ to $+9.1$ points under the production floor form ($+9.2$ at the in-sample calibration point; $+11.2$ under the competing-risks form). Inside the binding interval, $+5.6$ is the value at the mid-grid anchor, and measured corrections to the floor and the accounting basis move that value in both directions within the interval (Section~\ref{sec:identification}). The estimated cohort-hazard form lands at a comparable aggregate level but does not corroborate it: its own coefficient's sign is not statistically distinguishable from zero, it carries no valid production-spec uncertainty, and it is excluded from headline figures (Appendix~\ref{sec:patha}). That constraint is also not evenly borne. The marginal's per-balance intensity tilts toward below-740-FICO ($1.23\times$--$1.26\times$) and above-80-LTV ($1.23\times$) borrowers (Section~\ref{sec:identification}), and the statutory assumability carve-out that would relieve it is bounded above by the 20.4\% Ginnie share, with take-up unmeasured; the mobility cost falls disproportionately on the households with the fewest refinancing alternatives.

How much of the remainder sits outside household-level decision-making is a question the cross-design test of Section~\ref{sec:robustness-crossdesign} leaves open, and its answer is asymmetric. The hazard framework's aggregate recovery survives a fully synthetic population, 106.0\% versus 107.0\%, which is evidence that its survival structure rather than its data source carries the aggregate level; its path-level timing does not transfer. The ABM's share moves with calibration choice and data source: 59.3\% recalibrated, a 60.2\% fifty-seed mean, 76.3\% at the book's composition. That range sits above the pre-committed threshold, and above it Section~\ref{sec:robustness-crossdesign} withdraws the paradigm reading---the contrast cannot, on that test alone, be cleanly attributed to modeling paradigm rather than to calibration and data source. On the paradigm question itself I therefore make no claim, and the path-level timing question remains open for every estimator. The decomposition this paper does deliver never rested on either.

U.S. policymakers weigh three considerations: (1) deep, homogenized secondary market liquidity (the TBA market); (2) continuous household mobility (avoiding the lock-in effect); and (3) convexity-neutral duration profiles (predictable institutional cash flows). This paper's evidence is that they do not form the strict trilemma a three-way ``trade-off'' would imply. That reading holds under its production floor form; the additive form roughly doubles the margin at issue. The U.S. has prioritized (1) at real cost to (2): household mobility is genuinely constrained, while the cash-flow shortfall against the caps is mostly rate-inelastic. The Federal Reserve staff's own record and the ex-ante projection corroborate that composition as composition, not as numerical agreement, and not independently of it, since the projection may already have embedded part of the channel (Section~\ref{sec:lit}).

What I no longer find support for is the assumption that (2) and (3) trade off sharply under face accounting, in which fixing mobility would cost face-value cash flow. Under the rule-only transplant, Appendix~\ref{sec:abm-danish}'s recalibrated Danish counterfactual finds that a market-value repurchase rule would have relieved (2), while its effect on (3) is incidence-conditional: an $\approx\$61$ billion improvement under face accounting, a $-\$51.0$ billion cost under cash accounting at the re-derived discount (the foregone household-to-bondholder transfer). The dissolution rests on thinner ground than the framing suggests, so I state it plainly. The rule-only Danish gap is the lock-in marginal viewed from the counterfactual side, so ``(2) and (3) do not trade off sharply'' reduces to ``the marginal is small''. Under the production floor form the marginal is bounded small by construction, a calibrated ceiling rather than an independent finding; under the disclosed additive form, which never censors the elasticity, the same counterfactual is about $+11$ points and nearly floor-invariant (Section~\ref{sec:robustness-floor}). What survives across forms is that the trade-off is modest, not that it is absent.

On either form, the binding constraint against adopting such a rule is not a mobility--cash-flow trade-off but secondary-market liquidity. The Danish system relies on strict one-to-one match funding (the Balance Principle); the United States relies on the To-Be-Announced (TBA) forward market, which requires highly homogenized, pooled mortgages, and \citet{vickery2013} show that this homogeneity produces a measurable liquidity premium, with TBA-eligible pools trading at significantly lower cost than comparable ineligible specified pools. What a market-value repurchase feature strains is not homogeneity itself but its \emph{locus}: Denmark standardizes at the level of large, tap-issued bond series, the TBA market at the level of the deliverable pool, and the Danish callable market's own depth shows that the two coexist under the series convention \citep{berg2018}. The binding cost is therefore migrating U.S. issuance architecture from pool-level to series-level standardization, not a liquidity impossibility, though adopting the rule without that redesign would risk a documented source of liquidity in the transition. That is an argument from the cited literature, not a result of this design: nothing here prices a series-level call, so the investor-side cost of the transplant is named and left unpriced (\citet{diep2021} show that prepayment risk is compensated in the cross section of MBS returns, which is the price this design does not measure).

As long as the U.S. maintains the par-payoff MBS structure within the TBA framework, severe duration extension should be expected to recur in future tightening cycles that begin, as this one did, with a large outstanding stock of deeply in-the-money, below-market-coupon mortgages. The state dependence documented by \citet{berger2021} and \citet{eichenbaum2022} makes the extension conditional on that stock, not on tightening per se, and the recurrence claim is itself an extrapolation beyond what this design tests. The accounting machinery is also episode-shaped: with essentially the whole book out of the money, the marginal's sign is forced and the refinancing margin is inert, so a mixed-moneyness window would require re-deriving both. This design tests one episode. The duration extension itself, however, is mostly the arithmetic of scheduled amortization and baseline turnover against caps set above any plausible prepayment path; an alternative payoff rule would have trimmed it by the lock-in marginal, not eliminated it, and would carry the liquidity cost above.

\subsection{Limitations and Avenues for Future Research}\label{sec:limitations}

Several limitations remain on both sides of the framework.

\textbf{Housing supply as a competing channel}: The mobility margin this paper imports is estimated in an environment of historically tight for-sale inventory, and housing supply enters here only as a friction scaler on the moving decision, never as a channel in its own right. A household that does not move because nothing is listed is observationally identical, in these data, to one that does not move because its coupon is below market; this design does not separate them, the direction is argued rather than measured, and a binding supply constraint would absorb some of what is attributed to lock-in. That is why the scaled-null exercise that removes the non-rate share of the moving decision is reported alongside the headline rather than as a robustness footnote (Section~\ref{sec:identification}; \citealp{aladangady2024,hsieh2019}); no inventory series is used to price it.

\textbf{Distributional silence}: This paper measures a book-wide institutional cost and never asks who bears it: Ginnie Mae borrowers---20.4\% of SOMA face, and the population whose structurally faster speeds Section~\ref{sec:pathb} documents---sit outside the Freddie Mac sample entirely, so nothing here establishes whether the mobility cost falls differently on FHA and VA households than on the conventional borrowers the sample contains. I record that as a distributional limitation of the design, not merely as a coverage bound on the estimate.

\textbf{Sample scale and comparability}: The ABM's 10,000-household synthetic population and the hazard framework's 75,000-loan real Freddie Mac sample differ in both size and data source; Section~\ref{sec:robustness-scale} settles the size question on the production specification, since population size changes precision, not central tendency. The completed cross-design pair addresses the data-source question asymmetrically: the hazard framework's aggregate recovery is largely data-source-independent, the ABM's is not (Sections~\ref{sec:robustness-crossdesign} and~\ref{sec:robustness-synthesis}).

\textbf{Data}: A dedicated ingest pipeline integrates the Freddie Mac 2017--2021 loan-level data (20 origination quarters) that both hazard paths use, and where real data files are unavailable, results are validated against a synthetic fixture before real-data estimation. Fifteen-year MBS, approximately 9\% of SOMA face value, remain excluded from cohort modeling on the hazard side only; the ABM side folds them in (Appendix~\ref{sec:robustness-pipeline}).

\textbf{Model}: On the ABM side, the monthly CPR path fit remains weak (raw $R^2$ of $-6.984$), despite a directionally sensible level response to rates, and the debt-to-income check remains front-end only, without a measure of total household debt. Appendix~\ref{sec:abm-danish}'s production Danish counterfactual is a rule-only transplant onto U.S. households that does not estimate how Danish institutional frictions would themselves migrate, and its Danish-level bracketing anchor shows the gap's sign is not robust to bundling baseline-mobility differences into the transplant. Its recalibration runs against \citepos{berger2026} separately-estimated elasticities, an unrefereed working paper, draft dated July 2026, that is moving: the general-equilibrium refinance-in-place estimate rose roughly twentyfold between its January and July 2026 drafts. On the hazard side, Path A's Poisson pseudo-maximum-likelihood estimator requires a mild ridge penalty to remain well-conditioned at 295 stratum fixed effects, a stabilization device rather than a specification choice, and a lag-0 CPR correlation of $-0.378$, more negative even than the ABM's $-0.318$, confines its evidential contribution to the aggregate dollar level (Appendix~\ref{sec:patha}). Path B's literature-calibrated hazard magnitudes are directionally and quantitatively closer to the empirical benchmark than any ABM specification, but they are simulated on the 75,000-loan real Freddie Mac sample with elasticities calibrated from secondary literature rather than estimated jointly with the empirical benchmark itself; its recovery percentage is reported across the calibration band (105.9\%--108.2\%), its full 5.5\%--7.7\% elasticity band is propagated through the microsimulation, and loan-level sampling uncertainty now enters from an ex-ante-specified stratified bootstrap over the 75,000-loan draw (Section~\ref{sec:pathb}). An external second-agency replication (Fannie Mae loan-level data, same pipeline end-to-end) reproduces the lock-in marginal at $+8.68$ points against Freddie's $+9.20$, inside its pre-committed envelope---a sign-and-envelope replication, not a hazard-level match, and one whose envelope gate is 11.06 points wide, so the 0.52-point agreement rather than the gate is what should be read from it.

\textbf{No outcome-holdout months}: No outcome-holdout months exist anywhere in this paper, one global fact about the design each of whose components is disclosed only locally. The nearest construction is the marginal's temporal holdout (Section~\ref{sec:robustness-floor}), where the baseline turnover floor is calibrated on the window's first nineteen months and the central-minus-null difference is summed over the remaining twenty-three alone---months that enter that calibration nowhere---returning $+\$45.1$ billion; that construction holds for the marginal, not for the level, and it is an input-stability check rather than an outcome validation, since realized prepayment from the held-out months enters neither simulated leg and so cannot falsify the marginal against data. The production baseline turnover floor is still calibrated on 2023--24 turnover inside the evaluation window, though the pre-committed off-window re-measurement finds the defensible off-window level higher than 4\%, so the direction of any anchoring error shrinks rather than inflates the identified marginal, and the headline marginal is now reported at that off-window level rather than the in-window one. Path A trains on 19 of the 42 QT months, and its ridge-selection holdout doubles as its out-of-sample diagnostic. Path B's \emph{level}, as distinct from its marginal, is still set by a floor measured in the window on which it is evaluated, so every recovery percentage in the paper is an in-sample or calibrated quantity in the months dimension, disciplined by ex-ante-specified sweeps, nulls, parity gates, and bootstraps rather than by clean holdout months.

\textbf{Identification}: The lock-in marginal, at either floor, is a model-based decomposition conditional on the disclosed hazard structure (Section~\ref{sec:identification}); causal identification in the econometric sense would require exogenous variation in rate gaps across otherwise-comparable loans, and this design does not exploit it. All headline figures are frozen in tagged run manifests, with reproduction details recorded in the replication package's run ledger.

\textbf{Composition and vintage}: The Ginnie composition bound has been upgraded from two published snapshots to a time-varying overlay built from the published monthly CPR series, which leaves the conventional-share marginal at $+7.3$ points at the in-sample calibration point ($+4.4$ composed with the off-window headline floor), positive under every variant, with the Ginnie-specific level contribution inside the static bound (Section~\ref{sec:pathb}). Replication alone could not retire the vintage residual: the 2022 and pre-2017 vintages (23.1\% and 10.6\% of book face) lie outside the 2017--2021 universe of both agencies' samples used here, and the observed-speed bound of Section~\ref{sec:pathb} now disciplines that extrapolation directly (\$11.7 billion, 1.5\% of benchmark, signed toward overstating trapped liquidity). What remains open is therefore no longer the absence of a bound but the bound's own proxy assumption---observed Fannie segment speeds standing in for the SOMA book's, the same single-agency posture the production calibration itself carries---since a modeled overlay in place of that proxy is not constructible under the committed sampler, and the 2022 leg of the observed-speed bound (\$3.3 billion of the \$11.7 billion) stands as that vintage's discipline. Appendix~\ref{app:composition} collects the formal composition comparison behind these bounds.

\textbf{Cyclical floor}: The cyclical-floor variant is complete, but what its grid measures is floor dispersion under the hard maximum rather than cyclicality, with the surviving co-movement component small and bounded (Section~\ref{sec:pathb}; Table~\ref{tab:uncertainty} notes; Appendix~\ref{sec:robustness-seasonalfloor}).

Three questions remain open. The first is whether to add a seasonal multiplier to Path B's baseline hazard (Section~\ref{sec:pathb}). The second is a total-debt measure, or threshold sensitivity, for the ABM's front-end DTI gate (Appendix~\ref{sec:method-frictions}). The third is the monthly-frequency question the timing diagnosis leaves: the empirical path's variation at that frequency is seasonal, not rate-driven, and no estimator here models it.

\clearpage
\bibliographystyle{plainnat}
\bibliography{references}

\clearpage
